
The concepts outlined in Chapter \ref{chap:chern_intro} -- Berry phase, curvature and Chern number -- were first proposed in the 1980s to explain the newly discovered quantum Hall effect \cite{klitzing_new_1980, laughlin_quantized_1981}, where they provided the mathematical underpinning for the TKNN invariant that defines the Hall conductance \cite{thouless_quantized_1982,berry_quantal_1984}. In the decades since, the Chern number has been one of the central tools for classifying condensed matter systems outside of the Landau symmetry-breaking framework \cite{haldane_model_1988, hasan_topological_2010,qi_topological_2011,chiu_classification_2016}. They have found use in describing ocean waves \cite{deplace_topological_2017} coupled mechanical systems \cite{nash_topological_2015,xiao_geometric_2015, mitchell_amorphous_2018}, as well as providing the foundations for the modern theory of polarisation \cite{resta_macroscopic_1999,vanderbilt_electric_1993, vanderbilt_berry_2018}. 

One of the defining characteristics of quantum Hall physics is the robustness of the edge states and Hall conductance to impurities in the sample, where certain amount of disorder is even \textit{required} to stabilise the quantised conductivity \cite{trugman_localization_1983,chalker_percolation_1988}. Despite this, our definition of the Chern number cannot be calculated in a disordered material. In order to calculate the Berry curvature -- and thus the Chern number -- we must work in momentum space. If a material lacks translational symmetry, it is not possible to apply Bloch's theorem, crystal momentum is not a good quantum number and we simply have no concept of momentum space. Disorder-resistant methods exist~\cite{niu_quantized_1985, bellissard_noncommutative_1994}, however they compute a global Chern number, and cannot capture the local properties of materials with compound structure. This tension, that our tools for describing disorder-resistant physics are themselves undermined by disorder, suggests that our mathematical framework is incomplete. 
 
One potential solution to this problem is the development of real space methods for calculating the Chern number. These attempt to translate the Chern number into a mathematical language that is resolved in real space, ensuring it is well-defined when the system has no translational symmetry. There are a litany of proposals, such as Kitaev's Chern marker \cite{kitaev_anyons_2006,mitchell_amorphous_2018}, the Bott Index \cite{toniolo_equivalence_2017,toniolo_time-dependent_2018,hastings_almost_2010, toniolo_bott_2022}, the Chern marker \cite{bianco_mapping_2011,bianco_orbital_2013}, the localizer index \cite{loring_guide_2019} and most recently structural spillage \cite{munoz-segovia_structural_2023}. In all of these cases, the aim is to re-express the definition of Chern number in terms of a set of real-space operators which are local, and thus do not rely on the notion of momentum space to make sense. Such a quantity should be equivalent to the Chern number when calculated in a translationally symmetric context; however would still be possible in cases where these conditions have been relaxed, and the notion of $k$-space is gone. These quantities afford the researcher tools with which to comment on the topological characteristics of materials that had previously been out of reach, such as amorphous crystals \cite{agarwala_topological_2017}, quasi-periodic lattices \cite{huang_theory_2018}, systems out of equilibrium \cite{caio_topological_2019,toniolo_time-dependent_2018,golovanova_truly_2021,markov_locality_2023} and even strongly interacting systems \cite{markov_local_2021}.

In this section we start by introducing three such quantities, the Kitaev Chern number, the Bott index and the Chern marker, giving their definitions as well as explaining how they are related.

\begin{shaded}
    add a bit more summary of wot u writing
\end{shaded}